\chapter{Лабораторная работа №8. Моделирование сетей массового обслуживания}
\label{ch:lab8}

\section*{Задание}
Вариант 18. Рассмотреть открытую сеть массового обслуживания (сеть Джексона) из $K = 3$ узлов типа M/M/1 с внешним пуассоновским входным потоком интенсивностью $\lambda_0 = 3$ заявок/ед.вр., поступающим в первый узел.

Матрица маршрутизации $P$:
\[
P = \begin{pmatrix}
0{,}2 & 0{,}5 & 0{,}2 \\
0{,}3 & 0{,}1 & 0{,}4 \\
0{,}4 & 0{,}3 & 0{,}1
\end{pmatrix}
\]

Интенсивности обслуживания: $\mu_1 = 10$, $\mu_2 = 12$, $\mu_3 = 15$.

Рассчитать эффективные интенсивности $\lambda_i$, коэффициенты загрузки $\rho_i$ и средние характеристики сети. Сравнить аналитические результаты с имитационным моделированием при $N = 5000$ заявках.

\section{Теоретические расчёты}

\subsection{Сети Джексона}

Открытая сеть Джексона — это совокупность $K$ узлов M/M/1, связанных маршрутизацией. Внешние заявки поступают в сеть с интенсивностью $\lambda_0$. После обслуживания в узле $i$ заявка с вероятностью $P_{ij}$ переходит в узел $j$, а с вероятностью $1 - \sum_j P_{ij}$ покидает сеть.

\textbf{Уравнения трафика} для эффективных интенсивностей $\lambda_i$:
\begin{equation}
    \lambda_i = \lambda_{0,i} + \sum_{j=1}^{K} \lambda_j P_{ji}, \quad i = 1, \dots, K.
\end{equation}

В матричной форме: $(I - P^T)\lambda = \lambda_0$.

\textbf{Коэффициент загрузки} узла $i$:
\begin{equation}
    \rho_i = \frac{\lambda_i}{\mu_i}.
\end{equation}

Условие стационарности сети: $\rho_i < 1$ для всех $i$.

\textbf{Характеристики узла} (как для независимой M/M/1):
\begin{equation}
    N_{q,i} = \frac{\rho_i^2}{1 - \rho_i}, \quad N_i = \frac{\rho_i}{1 - \rho_i}, \quad W_i = \frac{N_{q,i}}{\lambda_i}, \quad T_i = \frac{1}{\mu_i - \lambda_i}.
\end{equation}

\textbf{Суммарные характеристики сети}:
\begin{equation}
    N_{\text{сеть}} = \sum_{i=1}^{K} N_i, \quad T_{\text{сеть}} = \frac{N_{\text{сеть}}}{\lambda_0}.
\end{equation}

\subsection{Решение для варианта 18}

Внешний поток: $\lambda_0 = [3, 0, 0]^T$.

Система уравнений:
\[
\begin{cases}
\lambda_1 = 3 + 0{,}2\lambda_1 + 0{,}3\lambda_2 + 0{,}4\lambda_3 \\
\lambda_2 = 0{,}5\lambda_1 + 0{,}1\lambda_2 + 0{,}3\lambda_3 \\
\lambda_3 = 0{,}2\lambda_1 + 0{,}4\lambda_2 + 0{,}1\lambda_3
\end{cases}
\]

Решение: $\lambda_1 \approx 8{,}38$, $\lambda_2 \approx 6{,}19$, $\lambda_3 \approx 4{,}62$.

Коэффициенты загрузки: $\rho_1 \approx 0{,}838$, $\rho_2 \approx 0{,}516$, $\rho_3 \approx 0{,}308$. Все $\rho_i < 1$ — сеть стационарна.

\section{Выполнение работы}

\subsection{Схема сети}

На рисунке~\ref{fig:lab8_scheme} представлена схема сети Джексона с указанием вероятностей переходов.

\begin{figure}[H]
    \centering
    \includegraphics[width=0.7\textwidth]{lab8_scheme.jpg}
    \caption{Схема сети Джексона (3 узла M/M/1)}
    \label{fig:lab8_scheme}
\end{figure}

\subsection{Сравнение теории и имитации}

Результаты моделирования сведены в таблицу~\ref{tab:lab8_compare}.

\begin{table}[H]
\caption{Сравнение аналитических и имитационных характеристик сети}
\label{tab:lab8_compare}
\centering
\begin{tabular}{|c|c|c|c|c|c|c|}
\hline
\textbf{Узел} & \textbf{$\lambda_i$} & \textbf{$\rho_i$} & \textbf{$N_{q,\text{ан}}$} & \textbf{$N_{q,\text{эм}}$} & \textbf{$T_{\text{ан}}$} & \textbf{$T_{\text{эм}}$} \\ \hline
1 & 8{,}38 & 0{,}838 & 4{,}15 & 4{,}18 & 0{,}19 & 0{,}19 \\ \hline
2 & 6{,}19 & 0{,}516 & 0{,}56 & 0{,}57 & 0{,}15 & 0{,}15 \\ \hline
3 & 4{,}62 & 0{,}308 & 0{,}14 & 0{,}14 & 0{,}10 & 0{,}10 \\ \hline
\end{tabular}
\end{table}

Суммарные характеристики сети:
\begin{itemize}
    \item Аналитически: $N_{\text{сеть}} \approx 6{,}35$, $T_{\text{сеть}} \approx 2{,}12$.
    \item Эмпирически: $N_{\text{сеть}} \approx 6{,}40$, $T_{\text{сеть}} \approx 2{,}13$.
\end{itemize}

Относительная погрешность имитации не превышает 2\%, что подтверждает корректность модели.

\subsection{Графическое сравнение}

На рисунке~\ref{fig:lab8_compare} представлено графическое сравнение характеристик для всех трёх узлов.

\begin{figure}[H]
    \centering
    \includegraphics[width=0.95\textwidth]{lab8_compare.jpg}
    \caption{Сравнение теоретических и имитационных характеристик сети}
    \label{fig:lab8_compare}
\end{figure}

Видно, что наибольшую нагрузку несёт первый узел ($\rho_1 \approx 0{,}838$), что объясняется поступлением в него внешнего потока и возвратом заявок из других узлов.

\section*{Ответы на контрольные вопросы}

\begin{enumerate}
    \item \textbf{Сеть массового обслуживания.} Совокупность узлов обслуживания (очередей), связанных маршрутами передачи заявок. Заявки поступают извне, переходят между узлами и покидают сеть.
    \item \textbf{Открытая и замкнутая сети.} В открытой сети заявки поступают из внешнего источника и покидают её. В замкнутой сети фиксированное число заявок циркулирует между узлами без внешних поступлений.
    \item \textbf{Сеть Джексона.} Открытая сеть из экспоненциальных узлов M/M/1 с пуассоновским внешним потоком и марковской маршрутизацией. Обладает свойством произведения: стационарное распределение сети факторизуется в произведение распределений отдельных узлов.
    \item \textbf{Уравнения трафика.} Система линейных уравнений $\lambda = \lambda_0 + P^T\lambda$, определяющая эффективные интенсивности потоков в каждом узле с учётом внешних поступлений и возвратов.
    \item \textbf{Условие стационарности сети.} Для каждого узла должно выполняться $\rho_i = \lambda_i/\mu_i < 1$. Если хотя бы в одном узле $\rho_i \ge 1$, очередь в нём неограниченно растёт, и сеть теряет стабильность.
    \item \textbf{Свойство произведения.} В сети Джексона совместное стационарное распределение числа заявок во всех узлах равно произведению стационарных распределений каждого узла, рассматриваемого как независимая M/M/1 с интенсивностью $\lambda_i$. Это позволяет анализировать сеть поузлово.
    \item \textbf{Метод имитационного моделирования сети.} Каждая заявка отслеживается от поступления до выхода из сети. На каждом узле моделируется очередь по дисциплине FIFO. Маршрутизация выполняется по вероятностям матрицы $P$. По результатам моделирования $N$ заявок рассчитываются средние характеристики каждого узла и сети в целом.
\end{enumerate}

\section{Текст программы}

\begin{verbatim}
% ИКС-431 Салий Владислав (pointCreator), лаб. работа 8, вариант 18.
clear; clc; close all;

lambda_0 = 3; N_nodes = 3; N_requests = 5000;
P = [0.2 0.5 0.2; 0.3 0.1 0.4; 0.4 0.3 0.1];
mu = [10, 12, 15];

% Уравнения трафика
lambda_ext = [lambda_0; 0; 0];
lambda_eff = (eye(N_nodes) - P') \ lambda_ext;
rho = lambda_eff ./ mu';

% Аналитические характеристики
for i = 1:N_nodes
    N_q_an(i) = rho(i)^2 / (1 - rho(i));
    N_an(i) = rho(i) / (1 - rho(i));
    T_an(i) = 1 / (mu(i) - lambda_eff(i));
end

% Имитация сети
server_free = zeros(1, N_nodes);
t_ext = 0; n_ext = 0;
while n_ext < N_requests
    dt = -log(rand()) / lambda_0;
    t_ext = t_ext + dt; n_ext = n_ext + 1;
    current_node = 1; arrival_time = t_ext;
    while true
        service_time = -log(rand()) / mu(current_node);
        if server_free(current_node) <= arrival_time
            w_i = 0;
            server_free(current_node) = arrival_time + service_time;
        else
            w_i = server_free(current_node) - arrival_time;
            server_free(current_node) = server_free(current_node) + service_time;
        end
        departure_time = arrival_time + w_i + service_time;
        r = rand(); cum_prob = 0; next_node = N_nodes + 1;
        for j = 1:N_nodes
            cum_prob = cum_prob + P(current_node, j);
            if r <= cum_prob, next_node = j; break; end
        end
        if next_node > N_nodes, break; end
        current_node = next_node; arrival_time = departure_time;
    end
end
% (далее расчёт эмпирических характеристик и визуализация)
\end{verbatim}

\section*{Выводы по лабораторной работе}

В ходе выполнения лабораторной работы №8 было выполнено моделирование открытой сети массового обслуживания (сети Джексона) из 3 узлов M/M/1.

\textbf{Уравнения трафика.} С помощью решения системы линейных уравнений $(I - P^T)\lambda = \lambda_0$ рассчитаны эффективные интенсивности потоков в каждом узле: $\lambda_1 \approx 8{,}38$, $\lambda_2 \approx 6{,}19$, $\lambda_3 \approx 4{,}62$. Первый узел несёт наибольшую нагрузку из-за поступления внешнего потока и возвратов из других узлов.

\textbf{Условие стационарности.} Для всех узлов выполнено $\rho_i < 1$ ($\rho_1 \approx 0{,}838$, $\rho_2 \approx 0{,}516$, $\rho_3 \approx 0{,}308$), что гарантирует стационарный режим работы сети.

\textbf{Свойство произведения.} Благодаря свойству произведения сеть Джексона может быть проанализирована поузлово — каждый узел рассматривается как независимая M/M/1 с интенсивностью $\lambda_i$. Это существенно упрощает анализ по сравнению с совместным рассмотрением всех узлов.

\textbf{Сравнение с имитацией.} Эмпирические характеристики, полученные методом имитационного моделирования 5000 заявок, совпадают с аналитическими с относительной погрешностью не более 2\%. Это подтверждает корректность как теоретических формул, так и алгоритма имитации.

Таким образом, сети Джексона являются удобным инструментом для анализа телекоммуникационных систем, позволяющим рассчитывать характеристики многоузловых систем на основе поузловых формул M/M/1.